\documentclass[withoutpreface,bwprint]{cumcmthesis}

\usepackage{booktabs}
\usepackage{amsmath,amssymb}
\usepackage{graphicx}
\usepackage{xcolor}
\usepackage{float}

\title{题目名称}
\tihao{B}
\baominghao{xxxx}
\schoolname{XX大学}
\membera{}
\memberb{}
\memberc{}
\supervisor{}
\yearinput{2026}
\monthinput{09}
\dayinput{00}

\begin{document}
\maketitle

\begin{abstract}
本文针对……展开研究。由于……，直接采用……会导致……，因此从……出发建立统一模型框架。

\textbf{针对问题一，}……

\textbf{针对问题二，}……

\textbf{针对问题三，}……

\textbf{针对问题四，}……

\keywords{关键词1\quad 关键词2\quad 关键词3\quad 关键词4}
\end{abstract}

\section{问题重述}

\section{问题分析}

\section{模型假设}

\section{符号说明}

\section{模型的建立与求解}

\subsection{问题一}

\subsection{问题二}

\subsection{问题三}

\subsection{问题四}

\section{模型检验与稳健性分析}

\section{模型评价与推广}

\begin{thebibliography}{99}
\bibitem{ref1} 作者. 文献题名[J]. 期刊, 年, 卷(期): 页码.
\end{thebibliography}

\begin{appendices}
\section{程序与支撑材料说明}
\end{appendices}

\end{document}
